% ============================================================================
% 03_data.tex — Data
% ============================================================================

\section{Data}
\label{sec:data}

The analysis draws on official statistics for all housing inputs and on index-provider data for financial returns. All series are monthly, in nominal pounds sterling, and cover January 2005 to June 2026 (258 months). I focus on England and its nine regions: England is the only part of the United Kingdom for which official monthly rent statistics reach back to 2005.

The sample is bounded at both ends by data rather than by choice, and the two ends bind for different reasons. The start is set by rents: no regional private-rent series for England exists at monthly or quarterly frequency before January 2005, in any source, official or commercial. The end is set by prices: the UK House Price Index runs about two months in arrears, so June 2026 is the latest complete vintage even though rents are published to July. Appendix~\ref{app:data} provides construction details; Table~\ref{tab:sources} summarises the sources.

\begin{table}[H]
  \centering
  \caption{Data sources}
  \label{tab:sources}
  \begin{threeparttable}
  \begin{tabularx}{\textwidth}{@{}lXl@{}}
    \toprule
    \textbf{Series} & \textbf{Source} & \textbf{Coverage used} \\
    \midrule
    House prices & UK House Price Index, HM Land Registry (all dwellings, England and regions) & 2005:1--2026:6 \\
    Private rents & ONS Price Index of Private Rents (PIPR) and its historical series & 2005:1--2026:6 \\
    Mortgage rate & Bank of England quoted rates: five-year fix at 90\% loan-to-value, constructed (Section~\ref{sec:mortgage_rate}) & 2005:1--2026:6 \\
    Equity returns & MSCI ACWI net total return, GBP, less 0.12\% p.a.\ fund fee & 2005:1--2026:6 \\
    Gilt yield & 10-year UK government bond yield (OECD, via FRED) & 2005:1--2026:5 \\
    Margin rate & Bank of England Bank Rate plus 1.50\% (renter-leverage robustness) & 2005:1--2026:6 \\
    Stamp duty & HMRC SDLT schedules, first-time-buyer rules by month & 2003:12--2026:6 \\
    \bottomrule
  \end{tabularx}
  \end{threeparttable}
\end{table}

\subsection{House prices and rents}

House prices come from the UK House Price Index (UK HPI), the official transaction-based index constructed from HM Land Registry sales records by hedonic regression with geometric averaging \parencite{ukhpi}. I use the average price of all dwellings, monthly, for England and each region.

Rents come from the ONS Price Index of Private Rents (PIPR), the official measure of private rents based on Valuation Office Agency lettings data \parencite{ons_pipr,ons_bulletin2026}. Two features make PIPR the appropriate counterpart to the UK HPI. First, it publishes average rent \emph{levels} in pounds per calendar month, not only an index. Second, it is a \emph{stock} measure covering new and existing tenancies, matching the position of a long-term renter whose rent adjusts with the market rather than only at new lets. The live PIPR series begins in January 2015; earlier years use the ONS historical series, chain-linked to PIPR's predecessor index back to January 2005 for England and its regions \parencite{ons_pipr_hist}. The two sources agree exactly over their 122-month overlap, so the splice introduces no discontinuity.

\subsection{A composition-consistent representative dwelling}
\label{sec:representative}

Moving from regions to a single national figure raises a significant measurement issue. The headline UK HPI England price and the headline PIPR England rent each average the country, but they weight the regions differently: the price is a transaction-mix average, close to a simple average of the nine regions, whereas the rent is a tenancy-mix average dominated by South East and London's large private rental sector. In June 2026 the national average rent (\gbp{1{,}446}) exceeded the average rent of every region except London, while the national average price (\gbp{293{,}262}) sat near the middle of the regional range. The ratio of the two, a price-to-rent ratio of 16.9, consequently falls \emph{below the ratio of every individual region} (17.7 to 22.4). Any rent-versus-own comparison built directly on this pair prices a mid-England home against a more expensive London-weighted rent, overstating the cost of renting and biasing the comparison towards ownership.

I therefore construct the price and rent of a \emph{representative English dwelling} as the same population-weighted average of the nine regions:
\begin{equation}
  P^{\text{rep}}_t = \sum_{i=1}^{9} w_i\, P_{i,t}, \qquad
  R^{\text{rep}}_t = \sum_{i=1}^{9} w_i\, R_{i,t}, \qquad
  w_i = \frac{\text{pop}_i}{\sum_j \text{pop}_j},
  \label{eq:representative}
\end{equation}
with fixed weights from mid-2023 ONS regional population estimates. Both aggregates then share one regional composition, and the resulting price-to-rent ratio (20.7 in June 2026) lies within the regional range. Over the sample the representative dwelling appreciated from \gbp{162{,}312} to \gbp{319{,}488} (3.2\% per year), falling 18\% peak-to-trough during 2007--2009, while its rent rose from \gbp{692} to \gbp{1{,}286} per month (2.9\% per year). The regional analyses in Sections~\ref{sec:results} and~\ref{sec:rar} use each region's own price and rent and are unaffected by aggregation.

\subsection{Mortgage rates}
\label{sec:mortgage_rate}

The rate a first-time buyer pays is not the rate the average borrower pays. English first-time buyers overwhelmingly fix: of new first-time-buyer lending in the first half of 2025, 99.1\% was fixed-rate and 0.018\% originated on a reversion rate, and among the fixed loans the three-to-five-year band was the modal choice in six of the seven half-years the Financial Conduct Authority publishes \parencite{fca_psd}. The baseline borrower therefore takes a \textbf{five-year fix at 90\% loan-to-value}, refixed at the prevailing five-year rate every sixty months.

No single Bank of England series measures that: the five-year 90\% loan-to-value series begins in February 2019. Earlier months are constructed by interpolating in loan-to-value between the published five-year series at 75\% and 95\%, with a sixty-one-month gap during 2008--2013, when fewer than three lenders offered a 95\% five-year fix, filled from the two-year 90\% series. Appendix~\ref{app:data} gives the construction and its validation. The series averages 4.74\% and ranges from 2.33\% to 6.99\%.

Two features of it matter for the results. It sits 1.19 percentage points above the Bank's effective rate on all new advances (CFMBJ95). That series is volume-weighted across every loan-to-value tier and fixation length, so recasting it monthly removes the fixation cliff through which a British first-time buyer is actually exposed to rate moves; Section~\ref{sec:robustness} reports it as a bounding case, and the difference is large. And because lenders price in bands and a borrower moves between them, the rate is repriced at each refix by the borrower's \emph{current} loan-to-value. A buyer entering at 90\% on a thirty-year term is at 74.8\% by their first refix, and charging them the 90\% premium then, a premium that peaked at 2.2 percentage points over 2010--2014, would overstate the cost of owning by several percentage points in exactly the period that decides the comparison.

\subsection{Financial returns}

The renter's portfolio tracks the MSCI ACWI index, covering large- and mid-capitalisation equities in 47 developed and emerging markets \parencite{msci_acwi}. I use end-of-month index levels computed by MSCI directly in pounds sterling in the net-total-return variant (dividends reinvested after withholding taxes), less an ongoing charge of 0.12\% per year representing a low-cost global index fund available to UK retail investors. Measuring the index in sterling embeds the currency exposure an unhedged UK investor bears. Over the sample the net-of-fee series returned 10.4\% per year, a factor of 8.38. As a home-bias robustness check I also use a FTSE~100 total-return series, which returned 7.1\% per year, a factor of 4.38; the choice between the two turns out to be one of the most consequential in the paper, and Section~\ref{sec:robustness} reports it in full. The 10-year gilt yield (OECD, via FRED) serves as a conservative alternative opportunity cost in Section~\ref{sec:rar}; it ranged from 5.4\% (June 2007) to 0.2\% (July 2020) and stood at 4.9\% in May 2026, its latest observation \parencite{fred_uk_gilts}.

Figure~\ref{fig:market_overview} shows the model's housing inputs together with the price-to-rent ratio they imply; Figure~\ref{fig:stock_returns} compares cumulative asset performance. The gap the buyer must overcome is evident: over twenty-one years the global equity portfolio multiplied roughly eightfold, UK large-cap equities fourfold, and the representative English dwelling, before the owner's running costs, twofold.

\begin{figure}[H]
  \centering
  \includegraphics[width=\textwidth]{fig1_market}
  \caption{Housing-market inputs for the representative English dwelling, 2005--2026. Panel A: average house price. Panel B: average private rent. Panel C: price-to-rent ratio (price divided by annual rent); the dashed line marks the sample mean of 21.9. Panel D: Bank of England effective interest rate on new mortgage advances to households. Sources: HM Land Registry UK HPI; ONS PIPR; Bank of England.}
  \label{fig:market_overview}
\end{figure}

\begin{figure}[H]
  \centering
  \includegraphics[width=\textwidth]{fig2_assets}
  \caption{Cumulative performance of global equities (MSCI ACWI, net total return in GBP, less a 0.12\% fund fee), UK equities (FTSE 100 total return) and representative English house prices, January 2005 = 100. Sources: MSCI; Morningstar; HM Land Registry.}
  \label{fig:stock_returns}
\end{figure}
